\begin{textsample}{POS Dim 2 – human – Score 7.00 – rj/rj\_ef\_ciencias\_da\_natureza\_1\_p2.txt}  \label{ex:f2_pos_005}
O ensino de Ciências é um componente fundamental na formação do cidadão contemporâneo.
Segundo Chassot ( 2000 ), a Ciência é considerada “ como uma linguagem para facilitar a leitura do mundo ” e aspira, para além da leitura, a intervenção. “ As \textbf{necessidades} de transformá -lo – e, preferencialmente, transformá -lo em algo melhor ” ( CHASSOT, 2003 ).
A formação de um indivíduo crítico que \textbf{saiba} se posicionar ativamente e intervir na \textbf{sociedade} moderna depende de conhecimentos diversos, tais como éticos-políticos, culturais, filosóficos e sobretudo, científicos.
Sendo assim, ao longo do Ensino Fundamental, o ensino de Ciências tem um compromisso com o letramento científico. É o letramento científico, um importante exercício da cidadania que facilita a leitura, interpretação e compreensão do mundo e, também, de transformá -lo com base nos aportes teóricos e processuais das ciências.
Sendo assim, o currículo de Ciências deve promover situações que garantam aos estudantes o \textbf{desenvolvimento} das oito competências específicas para o ensino de Ciências da natureza ( articuladas com as dez competências gerais da Base ) descritas na BNCC ( 2017 ).

% matched lemmas: desenvolvimento, necessidade, saber, sociedade
\end{textsample}
